\section{Related Work}
\label{sec:related}

\paragraph{Workflow and agentic generation benchmarks.} Recent work
evaluates whether language models can generate multi-step workflows.
WorFBench~\cite{worfbench2025} benchmarks agentic workflow generation and
scores generated workflows largely at the level of graph structure and
node/edge correctness. Our work is complementary but distinct: we argue
that for a high-stakes execution domain such as DeFi, graph-level
correctness is necessary but insufficient, and we add executable and
on-chain safety layers on top of structural scoring.

\paragraph{Tool-use and API evaluation.} A related line evaluates LLMs as
tool users. ToolLLM~\cite{toolllm2024} and
API-Bank~\cite{apibank2023} measure whether models can select and invoke
external APIs to satisfy instructions. These benchmarks emphasize correct
API selection and invocation; they do not model domain-specific safety
predicates such as price-impact gating, nor do they execute the resulting
plans against a stateful financial system. We adopt the tool-use framing
but specialize the evaluation to safety-constrained on-chain execution.

\paragraph{Decentralized finance and automated market makers.} Our
execution layer builds on the constant-product automated market maker
design introduced by Uniswap~\cite{uniswapv2}, whose $x\cdot y = k$
invariant makes price impact an explicit, computable function of trade
size. This is precisely the property that lets us distinguish a slippage
bound from a price-impact gate and observe unsafe execution directly.

\paragraph{No-code DeFi workflow systems.} We evaluate an existing no-code
DeFi platform (Koan) as a first-class baseline rather than a strawman,
using its real intent parser and workflow generator. To our knowledge,
prior no-code DeFi tooling has not been evaluated against an executable
safety benchmark, which is the gap this paper addresses.
